%% Plan complexe : point image, module et argument (leçon p. 107-108).
%% M a pour affixe z = a + bi ; r = |z| = OM et theta = arg(z) est une mesure de l'angle
%% orienté (OI, OM). Le conjugué zbar a pour image le symétrique de M par rapport à l'axe réel.
\documentclass[tikz,border=4pt]{standalone}
\usepackage{liaisons}
\begin{document}
\begin{tikzpicture}[x=1.3cm,y=1.3cm,
  axe/.style={-{Stealth[length=5pt]}, line width=0.6pt},
  rappel/.style={line width=0.5pt, dash pattern=on 3pt off 2pt, black!55}]

\def\av{2}
\def\bv{1.5}

%% ---------------- repère ---------------------------------------------------
\draw[axe] (-0.9,0) -- (3.3,0);
\draw[axe] (0,-2.3) -- (0,2.5);
\node[font=\small, anchor=south east, inner sep=3pt] at (3.3,0.05) {axe réel};
\node[font=\small, anchor=east, align=right, inner sep=3pt] at (-0.12,2.0)
     {axe des\\imaginaires purs};
\node[font=\small, anchor=north east, inner sep=2pt] at (-0.04,-0.04) {$O$};
\fill (1,0) circle (1.5pt) node[font=\small, anchor=north west, inner sep=3pt] {$I$};
\fill (0,1) circle (1.5pt) node[font=\small, anchor=south east, inner sep=3pt] {$J$};

%% ---------------- traits de rappel -----------------------------------------
\draw[rappel] (\av,\bv) -- (0,\bv);
\draw[rappel] (\av,\bv) -- (\av,-\bv);
\draw[rappel] (0,-\bv) -- (\av,-\bv);
\node[font=\small, anchor=north, fill=white, inner sep=1.5pt] at (\av,-0.03) {$a$};
\node[font=\small, anchor=east, fill=white, inner sep=1.5pt] at (-0.05,\bv) {$b$};
\node[font=\small, anchor=east, fill=white, inner sep=1.5pt] at (-0.05,-\bv) {$-b$};

%% ---------------- vecteur OM, module et argument -----------------------------
\draw[-{Stealth[length=7pt]}, solideA, line width=1.4pt] (0,0) -- (\av,\bv)
     node[midway, sloped, above, font=\small, text=solideA, inner sep=3pt]
     {$r=\left|z\right|$};
\draw[solideA, line width=1pt] (0.75,0) arc (0:36.87:0.75);
\node[font=\small, text=solideA, anchor=west, inner sep=2pt] at (18:0.85) {$\theta=\arg(z)$};
\draw[black!55, line width=0.9pt] (0.52,0) arc (0:-36.87:0.52);
\node[font=\small, text=black!55, anchor=west, inner sep=2pt] at (-19:0.6) {$-\theta$};

%% ---------------- les deux points image --------------------------------------
\fill[solideA] (\av,\bv) circle (2pt);
\node[font=\small, anchor=south west, align=left, inner sep=3pt] at (\av,\bv)
     {$M(z)$\\$z=a+bi$};
\fill (\av,-\bv) circle (1.8pt);
\node[font=\small, anchor=north west, inner sep=3pt] at (\av,-\bv)
     {$M'(\overline{z})$};
\end{tikzpicture}
\end{document}
